\documentclass[aspectratio=169]{beamer}
\usetheme{Madrid}
\usepackage[utf8]{inputenc}
\usepackage[T1]{fontenc}
\usepackage{tikz}
\usepackage{amssymb}
\usepackage{graphicx}
\setbeamertemplate{navigation symbols}{}
\setbeamercolor{title}{fg=white,bg=teal!60!black}
\setbeamercolor{frametitle}{fg=white,bg=teal!60!black}
\setbeamercolor{block title}{bg=teal!60!black!20,fg=black}
\setbeamercolor{block body}{bg=teal!60!black!8,fg=black}
\setbeamercolor{item}{fg=teal!60!black}
\setbeamercolor{itemize item}{fg=teal!60!black}
\setbeamertemplate{itemize item}[triangle]

\title{Computer Basics 1 — What Is a Computer?}
\subtitle{CPU, RAM, storage, input, output — no code yet, just the real parts}
\author{ByteForge --- Grade 7 --- Game Maker}
\date{}

\begin{document}

\begin{frame}
    \titlepage
\end{frame}

\begin{frame}{Today's Plan}
    \begin{itemize}
        \item Match the Parts
        \item Fill In: CPU, RAM, Input
        \item Fill In: Output \& Storage
        \item[$\bigstar$] BOSS --- Mixed Test
    \end{itemize}
\end{frame}

\begin{frame}{The CPU is the "brain" — it carries out instructions}
    \large
    \begin{itemize}
        \item Every real computer has a CPU doing this exact job, whether it's a phone, a laptop, or a games console
        \item The part that actually executes instructions, one after another.
    \end{itemize}
\end{frame}


\begin{frame}{RAM is temporary — it clears when power turns off}
    \large
    \begin{itemize}
        \item RAM holds whatever is running RIGHT NOW; turn the power off and it's genuinely gone. This is what makes an unsaved document actually disappear after a crash.
    \end{itemize}
\end{frame}


\begin{frame}{Storage keeps files even with the power off}
    \large
    \begin{itemize}
        \item Unlike RAM, storage (a hard drive, an SSD) is genuinely still there after a full shutdown
        \item It's why your files are still there the next time you turn a computer on.
    \end{itemize}
\end{frame}


\begin{frame}{Input devices send information IN; output devices send it OUT}
    \large
    \begin{itemize}
        \item A keyboard and mouse send information TO the computer; a monitor sends information FROM the computer back to you
        \item A real, simple, universal split every device fits into.
    \end{itemize}
\end{frame}

\begin{frame}{Real-World Examples}
    \large
    \begin{itemize}
        \item Every real computer, phone, or games console genuinely has all five of these parts, even when they're hidden inside a sealed case
        \item A phone's CPU, RAM, and storage are just much smaller than a desktop's.
    \end{itemize}
\end{frame}


\begin{frame}{Let Them Discover}
    \Large
    Look at your own phone or a family computer. Can you name which part is doing the CPU's job, which is RAM, and which is storage — even though you can't actually see them?
\end{frame}

\begin{frame}{Your Turn}
    \Large
    Open the student map, find \textbf{Computer Basics 1 — What Is a Computer?}, and work through it
    hands-on --- everything you just saw, they now do themselves.
\end{frame}

\end{document}
